\thispagestyle{empty}
\section*{Kurzdarstellung}
\label{sec:kurzdarstellung}

Diese Arbeit untersucht, wie sich latente Räume in generativen Diffusionsmodellen systematisch erkunden und navigieren lassen. Obwohl latente Räume eine zentrale Grundlage moderner Bildgenerierung bilden, sind ihre innere Struktur und ihre praktischen Grenzen weiterhin nur schwer nachvollziehbar. Im Mittelpunkt der Untersuchung stehen Navigationsverfahren wie Interpolation und die kontrollierte Bewegung entlang ausgewählter Richtungen, ohne dabei das zugrunde liegende Modell selbst zu verändern.

Zur Unterstützung dieser Untersuchung wurde ein prototypisches Explorationswerkzeug als externe Erweiterung für ComfyUI entwickelt. Das Plugin ermöglicht die Erstellung promptbasierter Datensätze, die Projektion und Clusterung ihrer Embeddings sowie die Ausführung verschiedener Navigationsoperationen innerhalb eines reproduzierbaren Arbeitsablaufs. Dadurch entsteht eine praktische Umgebung, in der sich das Verhalten unterschiedlicher Navigationsverfahren vergleichen und ihre Ergebnisse visuell untersuchen lassen.

Der experimentelle Teil der Arbeit analysiert generierte Bildsequenzen und Navigationsergebnisse im Hinblick auf Kontinuität, semantische Veränderung und Stabilität. Die Ergebnisse zeigen, wie ein werkzeuggestützter Ansatz dazu beitragen kann, das Verhalten latenter Räume zugänglicher zu machen, und verdeutlichen zugleich die Grenzen der Interpretierbarkeit und Steuerbarkeit diffusionsbasierter generativer Systeme.


\section*{Abstract}
\label{sec:abstract}

This thesis examines how latent spaces in generative diffusion models can be explored and navigated. Latent spaces are central to modern image generation, but their internal structure and practical limits are still difficult to understand. The work focuses on navigation methods such as interpolation and controlled movement along selected directions without changing the underlying model itself.

To support this analysis, a prototype exploration tool was developed as an external extension for ComfyUI. The plugin makes it possible to create prompt-based datasets, project and cluster their embeddings, and run different navigation operations in a reproducible workflow. This provides a practical environment for comparing how different navigation methods behave and how their results can be inspected visually.

The experimental part of the thesis studies generated image sequences and navigation outputs with regard to continuity, semantic change and stability. The thesis shows how such a tool-based workflow can help make latent-space behavior more accessible, while also highlighting the limits of interpretability and controllability in diffusion-based generative systems.
